% !PRINTED 258
\weberpage{258}
%<<CONT

\[
f(\omega) + \frac{\sqrt{2}}{f(\omega)}
\]

depend. They are \gesperrt{polynomials} in this combination
if the $\Phi_{a,c,\partial}$ remain finite for every finite, non-real $\omega$,
and they are \gesperrt{constant} if all the $\Phi_{a,c,\partial}$ remain finite for $q = 0$ as well, that is,
for $f(\omega) = 0$ and $f(\omega) = \infty$. In this case
\gesperrt{all the $\Phi_{a,c,\partial}$ are equal to one and the same constant}
(\S\,70).

From this the theorem follows:

\gesperrt{If one forms polynomials $\Phi_{a,c,\partial}$ in}

\[
f\left(\frac{c+\partial\omega}{a}\right), \quad f(\omega), \quad \frac{1}{f\left(\dfrac{c+\partial\omega}{a}\right)}, \quad \frac{1}{f(\omega)},
\]

\gesperrt{which have the property:}

1. \gesperrt{by the substitutions}

\[
\left(\omega, -\frac{1}{\omega}\right), \quad (\omega, \omega+2), \quad \left(\omega, \frac{\omega-1}{\omega+1}\right)
\]

\gesperrt{of being only permuted among one another,}

2. \gesperrt{of not becoming infinite for $q = 0$, then}

\[
\Phi_{a,c,\partial} = \text{constans}
\]

\gesperrt{is a transformation equation.}

In order to satisfy these conditions, the influence
of the substitutions (1) upon the functions $f$ must first be investigated.
This follows from the compositions \S\,69, (8) to (13)
and from the transformation formulae for the $f$-functions \S\,40.

For brevity we introduce the notation:

\weq{3}{\begin{aligned}
f(\omega) &= u, & f\left(\frac{c+\partial\omega}{a}\right) &= v, \\[4pt]
f_1(\omega) &= u_1, & \left(\frac{2}{a}\right) f_1\left(\frac{c+\partial\omega}{a}\right) &= v_1, \\[4pt]
f_2(\omega) &= u_2, & \left(\frac{2}{\partial}\right) f_2\left(\frac{c+\partial\omega}{a}\right) &= v_2.
\end{aligned}}

There then result the following corresponding changes, if in the
notation no account is taken of the conversion of the
numbers $a, c, \partial$ into $a, c_1, \partial$ or $a_2, c_2, \partial_2$, as it has just been
characterised by the cited formulae of \S\,69:
%CONT>>
